%%%%%%%%%%%%%%%%%%%%%%%%%%%%%%%%%%%%%%%%%%%%%%%%%%%%%%%%%%%%
%%% ELIFE ARTICLE TEMPLATE
%%%%%%%%%%%%%%%%%%%%%%%%%%%%%%%%%%%%%%%%%%%%%%%%%%%%%%%%%%%%
%%% PREAMBLE 
\documentclass[9pt]{elife}
% Use the onehalfspacing option for 1.5 line spacing
% Use the doublespacing option for 2.0 line spacing
% Please note that these options may affect formatting.
% Additionally, the use of the \newcommand function should be limited.

\usepackage{lipsum} % Required to insert dummy text
\usepackage[version=4]{mhchem}
\usepackage{siunitx}
\usepackage{tikz}
\usetikzlibrary{arrows.meta, positioning, shapes, calc}
\usepackage{rotating}
\usepackage{changepage} % in preamble


\DeclareSIUnit\Molar{M}

%%%%%%%%%%%%%%%%%%%%%%%%%%%%%%%%%%%%%%%%%%%%%%%%%%%%%%%%%%%%
%%% ARTICLE SETUP
%%%%%%%%%%%%%%%%%%%%%%%%%%%%%%%%%%%%%%%%%%%%%%%%%%%%%%%%%%%%
\title{Artificial life: virtual zebrafish based on hiearchical active inference}

\author[1,2,3*]{Hae-Jeong Park}
\affil[1]{Graduate School of Medical Science, Brain Korea 21 Project, Department of Nuclear Medicine, Psychiatry, Yonsei University College of Medicine, Seoul, Republic of Korea}
\affil[2]{Center for Systems and Translational Brain Science, Institute of Human Complexity and Systems Science, Yonsei University, Seoul, Republic of Korea}
\affil[3]{Department of Cognitive Science, Yonsei University, Seoul, Republic of Korea}

\corr{parkhj@yonsei.ac.kr}{PHJ}


% \presentadd[\authfn{5}]{eLife Sciences editorial Office, eLife Sciences, Cambridge, United Kingdom}

%%%%%%%%%%%%%%%%%%%%%%%%%%%%%%%%%%%%%%%%%%%%%%%%%%%%%%%%%%%%
%%% ARTICLE START
%%%%%%%%%%%%%%%%%%%%%%%%%%%%%%%%%%%%%%%%%%%%%%%%%%%%%%%%%%%%

\begin{document}

\maketitle

\begin{abstract}
Larval zebrafish provide a rare combination of a near-complete whole-brain connectome, optical access to neural activity, and rich visually guided behaviors.
Here we introduce a large-scale spiking neural network constrained by the zebrafish structural connectome that learns and acts using only local information.
Each model neuron is endowed with two effective compartments---a dendritic prediction compartment and a somatic spike generator---and synapses adapt according to a somato-dendritic mismatch rule inspired by the neuronal least-action principle.
This rule implements online predictive learning in the recurrent circuit without backpropagation-through-time.
We embed the network in a minimal virtual environment comprising prey-like and predator-like stimuli and define a small discrete latent state space (food, enemy, neutral) and action repertoire (approach, flee, indifferent).
Beliefs about latent causes are read out from higher-order populations and used to evaluate policies under an expected free energy objective, providing an active inference interpretation of orienting and escape.
In closed-loop simulations, the connectome-based spiking network learns to discriminate visual contexts and generates appropriate tail motor biases using only local plasticity and a simple preference structure.
Our results suggest that biologically plausible dendritic learning rules, when combined with the zebrafish connectome, are sufficient to implement core elements of predictive processing and active inference in a realistic whole-brain model.

\end{abstract}



\section{Hierarchical Vision Architecture (Retina → Thalamus → Cortex → Tectum)}

The PC-SNN agent implements a multi-stage visual hierarchy closely inspired by
biological visual processing in zebrafish.  
This hierarchy transforms raw sensory observations into latent cortical features,
computes cross-modal novelty and salience, and finally generates visuomotor drive
through the optic tectum.  
The system consists of four computational layers:

\[
\text{RetinaPC} \;\rightarrow\; \text{ThalamusRelay} \;\rightarrow\; 
\text{VisualCortexPC} \;\rightarrow\; \text{OpticTectum}.
\]

Each stage performs predictive-coding computation, producing both latent
representations and prediction errors used throughout the active-inference loop.

% ---------------------------------------------------------
\subsection{RetinaPC: Sensory Prediction and Feature Encoding}

The retinal layer implements a temporal predictive-coding loop that transforms
incoming luminance and contrast signals into a compact feature representation.
Let $x_t$ denote the retinal input at time $t$, and let $W$ be the retinal encoding
matrix.  The predicted feature vector is

\begin{equation}
\hat{x}_t = x_t W,
\end{equation}

and the temporal prediction error is

\begin{equation}
\varepsilon_t = \hat{x}_t - \hat{x}_{t-1},
\end{equation}

yielding retinal free energy

\begin{equation}
F^{(r)}_t = \frac{1}{2}\|\varepsilon_t\|^2.
\end{equation}

Hebbian updates, optionally modulated by dopaminergic tone, refine the retinal
weights to minimize future surprise.  
The retina thus provides both a bottom–up feature code and a fast estimate of
visual prediction error.

% ---------------------------------------------------------
\subsection{ThalamusRelay: Cross-Modal Salience and Novelty}

The thalamic relay receives free-energy estimates from visual ($F_v$), auditory ($F_a$),
and optionally body-state ($F_b$) channels.  
Its role is to compute cross-modal salience (CMS), reflecting the degree of mismatch
or novelty across modalities.  
Temporal novelty is measured as

\begin{equation}
\Delta F_i(t) = |F_i(t) - F_i(t-1)|,
\quad i \in \{v,a,b\}.
\end{equation}

Cross-modal divergence is defined as

\begin{equation}
C = 
\begin{cases}
|F_v - F_a|, 
& \text{(dual mode)} \\[6pt]
\frac{|F_v - F_a| + |F_v - F_b| + |F_a - F_b|}{3}, 
& \text{(tri mode)}.
\end{cases}
\end{equation}

Raw salience integrates both components:

\begin{equation}
S_t = \gamma_{\Delta} \cdot \mathrm{mean}(\Delta F_i)
    + \gamma_{C} \cdot C,
\end{equation}

and is smoothed with multiplicative stochastic bursts:

\begin{equation}
\mathrm{CMS}_t
= (1-\alpha)\,\mathrm{CMS}_{t-1}
+ \alpha S_t (1 + \beta\xi_t),
\quad \xi_t \sim \mathcal{N}(0,1).
\end{equation}

CMS provides a novelty-aware signal to dopaminergic systems and to higher-order
cortical inference.

% ---------------------------------------------------------
\subsection{VisualCortexPC: Context- and Goal-Modulated Predictive Coding}

Above the retinal and thalamic layers, the visual cortex constructs a latent visual
representation that incorporates both sensory evidence and top–down contextual or
goal information.

Let $x$ be the fused or competitive hemispheric visual input.  
Encoding and reconstruction follow predictive-coding dynamics:

\begin{align}
z &= x W_{\mathrm{enc}}, \\
\hat{x} &= z W_{\mathrm{dec}}, \\
\varepsilon^{(c)} &= x - \hat{x}.
\end{align}

Working-memory context $c_t$ modulates the latent state:

\begin{equation}
z \leftarrow z + \lambda_{\mathrm{ctx}} c_t,
\end{equation}

while goal-directed feedback $g_t$ provides task-conditioned modulation:

\begin{equation}
z \leftarrow z + \lambda_{\mathrm{goal}} g_t.
\end{equation}

Hebbian updates refine both encoder and decoder weights, ensuring that cortical
representation remains sensitive to task structure, temporal continuity, and
environmental context.

% ---------------------------------------------------------
\subsection{OpticTectum: Visuomotor Policy Field}

The optic tectum transforms bilateral salience and cortical latent features into
oculomotor drive for pursuit and orienting.  
Let $v_L$ and $v_R$ denote left and right visual salience, and $\bar{F}_t$ the
average free energy.  
The tectal motor drive is

\begin{equation}
d_t = \tanh\!\big(1.5(v_R - v_L)\big)
      - \lambda_F \bar{F}_t
      + g_{\mathrm{BG}}\, b_t,
\end{equation}

where $b_t$ is basal-ganglia gating.  
Eye velocity is then updated with momentum:

\begin{equation}
\dot{x}_t = 0.85\,\dot{x}_{t-1} + 0.15\, d_t,
\end{equation}

and dopamine modulates the centering force

\begin{equation}
x_t = x_{t-1} + \dot{x}_t
      - \big(c_0 + \Delta c(D_t)\big) x_{t-1},
\end{equation}

with bounds preventing saturation.  
The tectum acts as the final visuomotor decision module, selecting pursuit
direction and magnitude to minimize prediction error.

% ---------------------------------------------------------
\subsection{Functional Summary of the Vision Chain}

The full vision stack is organized hierarchically:
\begin{enumerate}
    \item \textbf{RetinaPC} extracts low-level visual features and computes temporal surprise.
    \item \textbf{ThalamusRelay} integrates cross-modal novelty into a scalar salience signal.
    \item \textbf{VisualCortexPC} builds context- and goal-sensitive latent visual representations.
    \item \textbf{OpticTectum} converts visual salience and cortical predictions into motor drive.
\end{enumerate}

In combination, these stages support the agent's ability to perform
prediction-error minimization across sensory, contextual, and task dimensions
while generating biologically realistic visuomotor behavior.


% ==============================================================
\section{From Predictive Coding to Active Inference (v1–v30)}
% ==============================================================

\subsection{Overview}

Versions v1–v30 represent the construction of a full embodied predictive-coding and active-inference architecture.
Beginning with unimodal sensory prediction (v1–v5), the model gradually incorporated cross-modal fusion (v6–v10),
motivational and dopaminergic learning (v11–v15),
goal selection and persistence (v16–v20),
interoceptive feedback and physiological regulation (v21–v25),
and finally temporal integration and metacognitive gain control (v26–v30).
Each version preserved backward compatibility through the \texttt{BaseModule} interface,
allowing seamless recombination of perception, action, and internal drives.

\subsection{v1–v5: Sensory Predictive Coding Core}

The early models implemented unimodal predictive-coding networks:
a \textbf{retinal encoder} and an \textbf{auditory predictive loop}.
Each sensory layer estimated latent causes $z_t$ from observations $o_t$
and minimized local free energy

\[
F_t = \frac{1}{2}\|o_t - \hat{o}_t\|^2 + \frac{1}{2}\|z_t - \hat{z}_t\|^2 ,
\]
via Hebbian weight updates
$\Delta W \propto o_t^\top z_t - \lambda W$.
RetinaPC modules handled spatial prediction errors,
while AudioPC modules captured temporal regularities.
These layers formed the perceptual foundation for later multimodal extensions.

\subsection{v6–v10: Multisensory Integration and Thalamic Novelty}

Visual and auditory streams were fused within a \textbf{ThalamusRelay} module that computed
cross-modal salience (CMS) as a function of temporal novelty and intermodal discrepancy:
\[
\mathrm{CMS}_t = \alpha \, (|\Delta F_v| + |\Delta F_a|)/2 + (1-\alpha)|F_v - F_a|.
\]
This signal modulated dopaminergic precision and triggered exploratory behavior.
Tri-modal extensions (v9–v10) added proprioceptive and body feedback terms,
enabling early forms of self–other distinction through multimodal prediction errors.

\subsection{v11–v15: Dopamine, Reward Prediction, and Policy Coupling}

The \textbf{DopamineSystem} generalized scalar reward prediction error (RPE)
into a precision-weighted motivational field:
\[
\delta_t = (r_t - \hat{r}_t) + \kappa \, \mathrm{CMS}_t .
\]
Coupled with the \textbf{BasalGanglia} module,
this created an alternation–exploration loop for saccadic or locomotor control.
Version v13 introduced \textbf{goal-dependent working memory},
allowing dopaminergic bursts to strengthen memory retention.
v14–v15 incorporated meta-gain adaptation,
so sustained dopamine increased consolidation time,
linking reward expectancy to memory persistence.

% ------------------------------------------------------------
\subsection{Hierarchical Architecture Summary}

\begin{verbatim}
[Environment]
   ↑ ↓
[Body / Interoception]  ←→  [Basal Ganglia / Motor]
   ↑ ↓
[Visual & Auditory Cortex] ←→ [Working Memory]
   ↑ ↓
[Dopamine System] ←→ [Policy Field / Active Inference]
   ↑ ↓
[Meta-Precision Monitor / Confidence Loop]
\end{verbatim}

This architecture unifies perception, action, emotion, and metacognition within a single active inference hierarchy, forming the computational substrate for a self-regulating, embodied cognitive agent.

% ------------------------------------------------------------
\subsection*{Footnote}

\footnotesize
\noindent
\textbf{Footnote.}~
The version numbering (v15–v30) follows a cumulative developmental trajectory of the
Predictive-Coding Spiking Neural Network (PC-SNN) framework.  
Each stage reflects the integration of additional inferential hierarchies:
from dopamine-based goal learning (v15) to active inference with expected free energy (v16),
temporal and social inference (v17–v19),
interoceptive embodiment (v20–v23),
adaptive environmental learning (v24–v26),
and finally meta-cognitive precision control and symbolic self-representation (v27–v30).
This systematic expansion parallels the evolution of biological nervous systems:
from sensory prediction and motor adaptation to the emergence of self-monitoring and conscious goal regulation.
\normalsize

---------

% ============================================================
\section{System Evolution: v15–v30 — From Goal Learning to Metacognitive Embodiment}
% ============================================================

The progressive versions (\textbf{v15–v30}) of the PC-SNN architecture represent a systematic expansion from sensory predictive coding and dopaminergic reinforcement to full hierarchical active inference and metacognitive embodiment. Each developmental stage integrates additional inferential depth—from simple reward prediction error (RPE) learning to self-monitoring of precision over precision—reflecting an increasingly realistic model of cognitive and affective brain function.

% ------------------------------------------------------------
\subsection{Version Map and Functional Expansion}

\begin{table*}[h!]
\centering
\caption{Evolution of Functional Architecture (v15–v30)}
\renewcommand{\arraystretch}{1.2}
\begin{tabular}{p{1.8cm} p{3.5cm} p{6.5cm} p{3.5cm} p{3.5cm}}
\hline
\textbf{Version} & \textbf{Core Innovation} & \textbf{Description} & \textbf{Biological Analogue} & \textbf{Expected Output} \\
\hline
v15.0–15.3 & Goal-Value Learning and Meta-Gain & Dopaminergic Q-value update and motivational gain stored in memory (WorkingMemory × PolicyField coupling). & VTA–PFC reinforcement loop & Goal persistence, value-dependent recall \\
v16.0–16.2 & Active Inference Policy Selection & Policies inferred via Expected Free Energy (EFE); integrates epistemic (uncertainty) and extrinsic (reward) terms. & ACC / medial PFC & Adaptive exploration vs. exploitation \\
v17.0–17.4 & Temporal Inference and Sequence Modeling & Adds short-term temporal filters (Kalman-like updates) and recurrent EFE tracking. & Dorsomedial PFC, premotor cortex & Predictive temporal planning \\
v18.0–18.9 & Social Valence Field (Multi-Agent Coupling) & Each agent’s dopamine field influenced by neighbors’ reward precision; introduces imitation and empathy weighting. & Superior temporal sulcus, mirror neuron systems & Alignment, cooperative or competitive dynamics \\
v19.0–19.5 & Joint Intent Inference (Social Active Inference) & Shared EFE minimization between agents, inferring both self and other’s beliefs; hierarchical coupling through precision channels. & TPJ / medial PFC & Theory-of-Mind-like inference \\
v20.0–20.3 & Temporal Active Inference (Interoceptive Loop) & Adds continuous EFE minimization with interoceptive prediction of internal states (e.g., energy). & Insula / brainstem homeostatic centers & Predictive homeostasis, energy regulation \\
v21.0–21.5 & Interoception + Metabolic Drive & Introduces glucose–O$_2$ metabolism variable that modulates dopamine and motivation. & Hypothalamus, orbitofrontal cortex & Hunger, fatigue, arousal modulation \\
v22.0–22.5 & Embodied Predictive Control (Body Model) & Adds musculoskeletal actuator dynamics (eye, tail, body torque) under predictive control. & Cerebellum / motor cortex & Coordinated saccade and locomotion \\
v23.0–23.9 & Hierarchical Action Control & Multi-timescale motor planning: reflexive (ms), rhythmic (100 ms), and goal-directed (s). & Motor hierarchy (SC–BG–PMC–PFC) & Smooth pursuit, goal-tracking trajectories \\
v24.0–24.5 & Heart and Respiratory Coupling (Cardiorespiratory AI) & Integrates heart rate variability and respiration into interoceptive free energy. & Brainstem / vagal–sympathetic loop & HRV-driven precision and fatigue modeling \\
v25.0–25.9 & Body–World Interaction (Environment Model) & Introduces dynamic environment (food, threat, light) → sensory likelihood updating. & Parietal cortex / hippocampus & Context-dependent goal pursuit \\
v26.0–26.5 & Learning under Uncertainty (Variational Policy Gradient) & Active inference implemented as a stochastic gradient on expected free energy with precision annealing. & ACC / OFC & Optimal uncertainty exploration \\
v27.0–27.5 & Hierarchical Intent and Meta-Cognitive Monitoring & Adds second-order inference on policy uncertainty; confidence controls dopamine release. & Dorsal ACC / frontopolar cortex & Meta-confidence oscillation \\
v28.0–28.9 & Social Reinforcement and Moral Valence & Collective reward and punishment; introduces “social EFE” weighting group coherence. & vmPFC / amygdala & Prosocial vs. antisocial trade-off \\
v29.0–29.5 & Language-Linked Symbolic Representation & Encodes symbolic cues (text/sound) as policies; forms early abstract concept model. & Broca / Wernicke areas & Symbol–goal alignment \\
v30.0+ & Full Embodied Conscious Loop & Integrates exteroception, interoception, proprioception, and meta-precision in unified self-model. & Whole-brain integration (DMN × SN × CEN) & Self-maintenance, metacognitive agency \\
\hline
\end{tabular}
\end{table*}

% ------------------------------------------------------------
\subsection{Stage I (v15–v16): Dopaminergic Goal Inference}

The system transitions from reactive dopamine-based reinforcement to expected free-energy (EFE)–based policy inference.  
Working memory becomes adaptive, adjusting retention by motivational gain, while the policy field evolves from simple Q-learning to probabilistic inference under uncertainty.  
At this stage, dopamine acts both as a prediction-error signal and a gain modulator influencing the temporal stability of memory traces.

% ------------------------------------------------------------
\subsection{Stage II (v17–v19): Temporal and Social Inference}

Temporal inference enables the agent to anticipate event sequences, while social coupling allows inference of others’ states.  
Each agent minimizes EFE jointly with neighbors, permitting imitation, empathy, and shared attention.  
This introduces the computational foundation for social valence pursuit and multi-agent active inference.

% ------------------------------------------------------------
\subsection{Stage III (v20–v23): Interoceptive Embodiment}

The generative model extends inward, including predictions of body states such as energy and proprioception.  
Interoceptive free-energy drives homeostatic regulation, while predictive control of eye and tail actuators allows adaptive locomotion.  
Here, action becomes perception through the minimization of proprioceptive prediction error.

% ------------------------------------------------------------
\subsection{Stage IV (v24–v26): Environmental Learning and Adaptive Precision}

Agents learn under uncertainty by estimating variational precision and minimizing expected free energy through stochastic gradient descent.  
Environmental contingencies (food, threat, light) are modeled explicitly, and adaptive annealing of precision enables robust exploration–exploitation balance.

% ------------------------------------------------------------
\subsection{Stage V (v27–v30): Meta-Cognition and Self-Model Formation}

Confidence dynamics (meta-precision) are added as a higher-order variable that monitors the reliability of inference itself.  
The agent thus develops an internal model of its own uncertainty.  
By v29–v30, the system integrates social, linguistic, and interoceptive layers into a unified self-model capable of maintaining goals, evaluating outcomes, and adjusting belief precision—approximating a computational model of metacognitive agency.

% ------------------------------------------------------------
\subsection{Equational Framework}

\begin{align}
\dot{s} &= -\frac{\partial F}{\partial s}, &
\dot{a} &= -\frac{\partial G}{\partial a} \\
F &= \text{Free Energy (perceptual inference)} \\
G &= \text{Expected Free Energy (policy inference)}
\end{align}

Each hierarchical level minimizes either perceptual free energy ($F$) or expected free energy ($G$), while the meta-cognitive level minimizes \textit{precision over precision}, adjusting how much confidence is assigned to inferential beliefs.


---------
\subsection{v16–v20: Goal Selection, Persistence, and Policy Learning}

These versions unified value learning and goal selection
within the \textbf{GoalPolicyField}.
Each goal $g_i$ maintained a Q-like value $Q_i$ updated by
temporal-difference error:
\[
Q_i \leftarrow Q_i + \eta (r_t + \gamma Q_i - Q_i).
\]
Policies were selected probabilistically
$\pi(g_i) \propto \exp(-\beta F_i)$,
combining precision, motivation, and novelty.
The working memory became goal-conditioned, receiving the one-hot goal vector
and maintaining contextual latent states over time.
Persistence timers stabilized chosen goals across several inference cycles,
creating a rudimentary notion of intent.

\subsection{v21–v25: Body, Interoception, and Homeostatic Control}

The system was extended with a \textbf{BodyModule}
representing muscles, viscera, and internal physiology.
Motor outputs (eye, tail, fin) were treated as proprioceptive predictions.
An interoceptive channel encoded heart rate, oxygen, and metabolic energy,
integrated by an \textbf{InteroceptivePC} subsystem.
Prediction errors in these internal signals
drove homeostatic regulation:
\[
F_{\mathrm{int}} = \sum_j \frac{(x_j - \hat{x}_j)^2}{2\sigma_j^2}.
\]
Coupling dopamine to metabolic free energy linked reward
to physiological cost, allowing adaptive trade-offs between
exploration and energy conservation.

\subsection{v26–v30: Temporal Integration and Meta-Gain Regulation}

Temporal inference was introduced through the \textbf{DeepTemporalField},
tracking changes in free energy across time and forecasting expected precision gains.
Each cycle updated a meta-gain variable
$g_t = 1 + \lambda(|\mathrm{RPE}_t| + |\mathrm{CMS}_t|)$,
which scaled both learning rate and memory retention coefficient.
High meta-gain corresponded to strong attention or arousal,
while low gain represented fatigue or disengagement.
This mechanism closed the perception–action–valuation loop,
producing coherent cycles of exploration, learning, and rest.

\subsection{Hierarchical Summary}

\begin{table}[h]
\centering
\caption{Major Functional Expansions (v1–v30)}
\begin{tabular}{l l l}
\hline
Version Range & Core Module Added & Functional Role \\
\hline
v1–v5 & RetinaPC, AudioPC & Sensory predictive coding \\
v6–v10 & ThalamusRelay & Cross-modal novelty and CMS gain \\
v11–v15 & DopamineSystem, BasalGanglia & Reward prediction, motivation \\
v16–v20 & GoalPolicyField, WorkingMemory & Goal inference and persistence \\
v21–v25 & Body/Interoceptive modules & Homeostasis and energy regulation \\
v26–v30 & DeepTemporalField & Temporal prediction and meta-gain control \\
\hline
\end{tabular}
\end{table}

\subsection{Summary}

Across v1–v30, the PC-SNN evolved from a reactive sensory predictor
into a unified active-inference agent.
Its architecture integrates hierarchical perception, multimodal fusion,
dopaminergic valuation, working-memory consolidation,
motor control, and interoceptive regulation.
Each subsystem participates in minimizing expected free energy:
perception reduces prediction error,
action fulfills sensory predictions,
and internal physiology maintains energy balance.
This establishes the foundation for higher-level social,
semantic, and metacognitive inference introduced in later versions (v31–v55).


% ==============================================================
\section{Unified Active Inference Architecture (v31.0–v33.0)}
% ==============================================================

\subsection{Overview}

The unified agent integrates all preceding modules (v1–v30) within a single
active inference framework. Each subsystem—vision, working memory,
dopaminergic policy control, interoceptive regulation, and meta-precision—
is formulated as a specialized free-energy minimization process.
The resulting architecture implements a multi-timescale
hierarchical generative model of perception, action, and self-regulation.

\subsection{Core Principle}

At every time step, the agent updates internal beliefs $\mu$,
actions $a$, and policies $\pi$ by minimizing a variational bound
on surprise, or \textit{free energy}:

\begin{equation}
F = \mathbb{E}_{q(s)}[\ln q(s) - \ln p(s,o,\pi)] ,
\end{equation}

where $s$ denotes latent states, $o$ sensory observations, and $\pi$
candidate policies.  Perception corresponds to optimizing posterior
beliefs $q(s)$; action corresponds to sampling sensations that fulfill
predictions; and policy selection minimizes the \textit{expected}
free energy $G(\pi)$:

\begin{equation}
G(\pi) = \mathbb{E}_{q(o,s|\pi)} \big[
D_{\mathrm{KL}}(q(s|\pi)\,||\,p(s))
 - \ln p(o|C)
\big],
\end{equation}

where $C$ encodes the agent’s prior preferences (goals).
This decomposition expresses the dual imperatives of
\textit{epistemic value} (reducing uncertainty) and
\textit{extrinsic value} (achieving preferred outcomes).

\subsection{Hierarchical Structure}

The unified agent comprises five nested layers, each minimizing
free energy on a distinct timescale:

\begin{itemize}
\item \textbf{L1: Sensory Predictive Coding}  
Visual and auditory cortices ($F_s$) infer latent causes of sensory input.
\item \textbf{L2: Working Memory Integration}  
Associative cortices ($F_m$) integrate short-term contextual information,
modulated by dopamine and CMS gain.
\item \textbf{L3: Policy and Motivation}  
Dopaminergic and striatal fields ($G_\pi$) implement goal selection via
expected free-energy minimization.
\item \textbf{L4: Interoception and Motor Control}  
Autonomic and somatic systems ($F_b$) regulate physiological and
kinematic states toward homeostasis.
\item \textbf{L5: Meta-Precision and Confidence}  
Prefrontal–DMN circuits ($F_\Pi$) monitor uncertainty and adapt belief
precision across levels.
\end{itemize}

Each level passes prediction errors upward and expectations downward,
forming a recursive active inference loop:

\[
\dot{\mu}_l = -\frac{\partial F_l}{\partial \mu_l}, \qquad
\dot{\pi} = -\frac{\partial G}{\partial \pi}, \qquad
\dot{a} = -\frac{\partial F}{\partial a}.
\]

\subsection{Computational Implementation}

All modules share a unified interface through the \texttt{BaseModule} class,
which maintains device management, learning rates, and diagnostic reporting.
The agent’s step function integrates seven sequential updates:

\begin{enumerate}
\item \textbf{Sensory Inference:}  
Predictive-coding loop in visual cortex minimizes $F_s$ given sensory inputs $o_t$.
\item \textbf{Contextual Memory:}  
Working-memory update $\mu_t = f(\mu_{t-1}, dopa, cms)$ modulates belief decay $\alpha_{\mathrm{eff}}$.
\item \textbf{Policy Inference:}  
GoalPolicyField computes $G(\pi)$ using dopaminergic RPE and memory state.
\item \textbf{Expected Free Energy:}  
ActivePolicyField evaluates epistemic and extrinsic components
$G_{\text{epistemic}}$ and $G_{\text{extrinsic}}$.
\item \textbf{Motor Actuation:}  
BodyMotorField samples an action $a_t$ from posterior $p(a_t|\pi_t)$.
\item \textbf{Interoceptive Regulation:}  
HeartInteroceptiveField adjusts internal energy and heart rate signals.
\item \textbf{Meta-Precision Update:}  
MetaPrecisionField updates global confidence $\Pi_t$ based on entropy of policy posteriors.
\end{enumerate}

All computations occur within a single GPU graph, with gradients detached
between hierarchical levels to emulate separate neural timescales.

\subsection{Mathematical Summary}

Each level contributes to the composite objective:

\[
\min_{a,\pi,\mu} \; F + G + H
\]

where:
\[
\begin{aligned}
F &\;=\; \text{Sensory and interoceptive free energy,} \\
G &\;=\; \text{Expected uncertainty (policy-level EFE),} \\
H &\;=\; \text{Entropy of meta-precision (confidence).}
\end{aligned}
\]

Thus, perception, learning, and metacognition are formally unified as
different gradients on the same variational bound.

\subsection{Validation Protocol}

To ensure continuity across all subsystems,
each component was subjected to modular testing.
Table~\ref{tab:validation_map} summarizes the key functional invariants.

\begin{table}[h]
\centering
\caption{Subsystem Validation under Active Inference Framework}
\begin{tabular}{l l l}
\hline
Subsystem & Optimized Quantity & Expected Behavior \\
\hline
Vision & $F_s$ (sensory free energy) & Prediction error $\downarrow$ \\
Memory & $\alpha_{\mathrm{eff}}$ (decay) & Retention $\uparrow$ with dopamine \\
Policy & $G_\pi$ (expected FE) & Goal selection consistency \\
Body & $F_b$ (motor FE) & Stable actuation \\
Heart & $F_i$ (interoceptive FE) & Homeostatic stability \\
Meta & $H$ (entropy of confidence) & Optimal precision regulation \\
\hline
\end{tabular}
\end{table}

\subsection{Summary}

The unified architecture demonstrates how sensory, motivational,
and metacognitive processes can all be expressed as hierarchical
inference under a single free-energy principle.
This framework provides a biologically interpretable and
computationally consistent foundation for subsequent
multi-agent and social active inference extensions (v34+).


\subsection{Unified Agent Testing Protocol (v31.0)}

To ensure functional continuity across all integrated modules (v1–v30),
we implemented a modular test harness where each subsystem is validated
independently under randomized synthetic inputs. Table~\ref{tab:test_map}
summarizes the test objectives, conditions, and expected invariant outcomes.

\begin{table}[h]
\centering
\caption{Subsystem Tests for Unified Agent v31.0}
\begin{tabular}{l l l}
\hline
Subsystem & Test Condition & Expected Invariant \\
\hline
Vision & Randomized input → Free energy > 0 & Predictive stability \\
Memory & High gain → Low α & Memory persistence \\
Policy & Dopamine ↑ → Choice bias ↑ & Goal modulation \\
Body & Random gain → Stable torque output & Smooth actuation \\
Heart & Energy deficit → HRV ↑ & Homeostatic correction \\
Meta & Confidence ↑ → Gain ↑ & Precision scaling \\
Environment & Action → Sensory feedback change & Closed loop consistency \\
\hline
\end{tabular}
\end{table}

% ==============================================================
\section{Temporal and Multi-Agent Active Inference (v34–v36)}
% ==============================================================

\subsection{Rationale}

Building upon the unified single-agent architecture (v31–v33),
versions v34–v36 extend active inference into the temporal and social domains.
Whereas previous models minimized instantaneous or short-term
free energy within an individual,
the new formulation introduces two key expansions:

\begin{enumerate}
\item \textbf{Temporal Active Inference:}  
Agents predict not only current states but the trajectories of future states,
minimizing the cumulative free energy across time.
\item \textbf{Multi-Agent Coupling:}  
Multiple agents share partially overlapping generative models,
allowing reciprocal inference about each other’s beliefs,
intentions, and actions.
\end{enumerate}

These extensions enable the modeling of collective behavior,
empathy, competition, and cooperation through shared precision-weighted
prediction errors.

\subsection{Temporal Generative Process}

Temporal active inference assumes a deep temporal generative model
$p(o_{1:T}, s_{1:T}, \pi)$ that unfolds across $T$ time steps.
The agent maintains beliefs over trajectories $\tilde{s} = (s_1, s_2, \dots, s_T)$
and selects policies that minimize the \textit{expected cumulative free energy}:

\begin{equation}
G(\pi) = \sum_{t=1}^{T} 
\mathbb{E}_{q(o_t, s_t|\pi)} 
\big[ \underbrace{D_{\mathrm{KL}}(q(s_t|\pi) || p(s_t))}_{\text{epistemic value}}
 - \underbrace{\ln p(o_t|C_t)}_{\text{extrinsic value}} \big].
\end{equation}

This formulation incorporates both the desire to reduce future uncertainty
and the drive to realize preferred sensory outcomes.
Temporal coherence is maintained through transition priors:

\[
p(s_t|s_{t-1}) = \mathcal{N}(f(s_{t-1}, a_{t-1}), \Sigma_t)
\]

that constrain plausible trajectories according to learned dynamics $f$.

\subsection{Temporal Implementation}

In v34, each agent maintains a bank of policy trajectories
$\pi_i = \{a_{t:t+H}^i\}$ over a fixed horizon $H$.
For each trajectory, expected free energy $G_i$ is computed and stored
in a rolling buffer:

\[
G_i = \mathbb{E}_{t:t+H} [ F_t + \lambda F_{t+1} + \lambda^2 F_{t+2} + \dots ],
\]
where $\lambda \in [0,1]$ is a temporal discount parameter.

The agent’s action selection probability is proportional to $\exp(-G_i)$,
yielding a Boltzmann distribution over temporal policies.
Policy persistence is achieved by a recurrent meta-field that
stores the expected precision of each trajectory, 
$\Pi_i = \exp(-\mathrm{Var}[G_i])$,
which controls exploration vs. exploitation.

\subsection{Multi-Agent Coupling}

Version v35 introduces multiple agents $\mathcal{A} = \{A_1, A_2, \dots, A_N\}$,
each minimizing its own free energy $F^{(n)}$ under partially shared latent causes.
For any agent pair $(A_i, A_j)$, the coupling term is defined as:

\begin{equation}
F_{\text{social}}^{(i,j)} 
= \eta \, \big\| \mathbb{E}_{q_i(s)}[s] - \mathbb{E}_{q_j(s)}[s] \big\|^2 ,
\end{equation}

where $\eta$ controls social precision (empathy gain).
This term penalizes divergence between inferred beliefs about
common latent states—such as shared goals, environmental affordances,
or perceived intentions.

Each agent thus optimizes:

\[
F_{\text{total}}^{(i)} = F^{(i)} + G^{(i)} + \sum_{j\neq i} F_{\text{social}}^{(i,j)} .
\]

This yields emergent cooperative, competitive, or alignment behaviors
depending on $\eta$ and prior preferences $C^{(i)}$.

\subsection{Collective Inference Loop}

At runtime, agents exchange precision-weighted beliefs
about shared latent causes via a communication field $C_t$:

\begin{align}
\dot{\mu}_i &= -\frac{\partial F^{(i)}}{\partial \mu_i}
              - \eta \sum_{j \neq i} (\mu_i - \mu_j), \\
\dot{a}_i   &= -\frac{\partial F^{(i)}}{\partial a_i}.
\end{align}

When $\eta$ is high, beliefs rapidly synchronize, yielding cooperation;
when $\eta$ is low or asymmetric, competition or deception can emerge.
This enables the simulation of social dynamics such as
leadership, following, altruism, or rivalry under a single formal principle.

\subsection{Meta-Precision and Collective Confidence}

Version v36 adds a collective meta-precision variable $\Pi_{\mathrm{group}}$
that tracks shared uncertainty across agents:

\[
\Pi_{\mathrm{group}} = \frac{1}{N} \sum_i \exp(-H[q_i(\pi)]),
\]

where $H[q_i(\pi)]$ is the entropy of each agent’s policy posterior.
Group confidence modulates collective gain and
determines whether agents synchronize or diversify their actions.
This provides a formal account of collective confidence,
crowd synchronization, and breakdown under ambiguity.

\subsection{Implementation Summary}

Table~\ref{tab:temporal_social_versions} summarizes the conceptual progression.

\begin{table}[h]
\centering
\caption{Temporal and Social Extensions of Active Inference Agent}
\begin{tabular}{l l l}
\hline
Version & Key Innovation & Functional Role \\
\hline
v34.0 & Temporal expected free energy & Future-oriented inference \\
v35.0 & Multi-agent coupling term & Social alignment / empathy \\
v36.0 & Group meta-precision & Collective confidence regulation \\
\hline
\end{tabular}
\end{table}

\subsection{Summary}

The temporal and multi-agent extensions generalize the unified agent
from an isolated predictive system to a distributed, temporally deep,
and socially coupled network.
Every process—anticipation, cooperation, competition, and
confidence alignment—emerges as a corollary of minimizing
expected free energy across agents and time.

% ==============================================================
\section{Physiological Active Inference (v37–v40)}
% ==============================================================

\subsection{Rationale}

Versions v37–v40 extend the active inference framework to include
autonomic and metabolic regulation.
In this extension, the body is not a passive actuator but an
integral part of the generative model:
internal organs, heart rate, respiration, and glucose metabolism
all participate in the same free-energy minimization process.
This formulation allows physiological homeostasis and cognitive control
to be expressed within a single variational principle.

\subsection{Interoceptive Generative Model}

Each agent maintains beliefs over internal physiological states
$s^{\mathrm{int}} = \{\text{heart rate}, \text{oxygenation}, \text{energy balance}\}$,
which evolve under stochastic dynamics:

\begin{equation}
p(s_t^{\mathrm{int}} | s_{t-1}^{\mathrm{int}}, a_t)
 = \mathcal{N}( f_{\mathrm{body}}(s_{t-1}^{\mathrm{int}}, a_t), \Sigma_{\mathrm{int}} ).
\end{equation}

Interoceptive sensory signals $o_t^{\mathrm{int}}$ (e.g., baroreceptors,
chemoreceptors, visceral afferents) are generated from these latent states:

\[
p(o_t^{\mathrm{int}}|s_t^{\mathrm{int}}) = \mathcal{N}( g_{\mathrm{body}}(s_t^{\mathrm{int}}), \Gamma_{\mathrm{int}} ).
\]

Prediction errors between observed and expected internal sensations drive
autonomic adjustments (heart rate, vasoconstriction, respiratory depth)
that minimize interoceptive free energy $F_i$.

\subsection{Homeostatic Set-Points as Prior Preferences}

Physiological stability is implemented as a prior preference distribution
$p(o_t^{\mathrm{int}}|C_i)$ centered around desired set-points $C_i$
(e.g., normal glucose, normoxia, 70–80 bpm).
Deviations from these priors generate negative valence signals
interpretable as discomfort, fatigue, or stress.

\begin{equation}
F_i = D_{\mathrm{KL}}(q(s_t^{\mathrm{int}})||p(s_t^{\mathrm{int}}))
     - \ln p(o_t^{\mathrm{int}}|C_i).
\end{equation}

This allows metabolic and affective variables to be unified as
prediction-error signals on internal priors.

\subsection{Autonomic Policy and Gain Control}

An autonomic policy field determines sympathetic vs.~parasympathetic
drive according to the gradient of interoceptive free energy:

\[
\dot{a}_{\mathrm{aut}} = -\frac{\partial F_i}{\partial a_{\mathrm{aut}}}.
\]

The resulting control signals modulate heart rate ($\mathrm{HR}$),
breathing amplitude ($\mathrm{BR}$), and vascular tone ($\mathrm{VT}$)
through efferent pathways.
The same dopaminergic precision mechanism used in cognitive levels
now regulates autonomic gain $\gamma_{\mathrm{aut}}$:

\[
\gamma_{\mathrm{aut}} = \gamma_0 \big( 1 + \kappa_{\mathrm{dopa}} (dopa - 0.5) \big).
\]

Thus, heightened dopaminergic activity increases metabolic readiness
and prolongs working-memory retention, linking motivation with energy supply.

\subsection{Energy Metabolism and Predictive Efficiency}

Version v38 introduces a metabolic field that tracks available energy $E_t$
and predicts its consumption under different policies $\pi$:

\[
E_{t+1} = E_t - \rho_{\mathrm{work}} a_t^2 + \rho_{\mathrm{rec}} s_t^{\mathrm{int}},
\]
where $\rho_{\mathrm{work}}$ and $\rho_{\mathrm{rec}}$ are expenditure
and recovery coefficients.

Energy becomes a precision variable: low energy increases sensory noise
and reduces policy confidence, producing fatigue-like behavior.
The agent thus learns to minimize a composite objective:

\[
F_{\mathrm{phys}} = F_i + \lambda_E \, |E_t - E_0|^2,
\]
which couples energetic balance with interoceptive stability.

\subsection{Heart–Brain Synchronization (v39)}

Cardiac oscillations are treated as a low-frequency rhythm
modulating cortical precision weights.
The phase $\phi_{\mathrm{HR}}(t)$ of the cardiac cycle entrains
slow oscillations in cortical gain parameters:

\[
\Pi_{\mathrm{ctx}}(t) = \Pi_0 \big( 1 + \epsilon \cos \phi_{\mathrm{HR}}(t) \big),
\]
producing cyclical windows of increased confidence and attention.
This mechanism accounts for physiological entrainment of cognition
and periodic fluctuations in predictive precision.

\subsection{Interoceptive–Cognitive Coupling (v40)}

In v40, the full agent integrates cognitive, emotional, and bodily layers
through reciprocal precision exchange:

\begin{align}
\dot{\mu}_{\mathrm{cog}} &= -\frac{\partial F}{\partial \mu_{\mathrm{cog}}}
                           + \omega \, (\mu_{\mathrm{int}} - \mu_{\mathrm{cog}}), \\
\dot{\mu}_{\mathrm{int}} &= -\frac{\partial F_i}{\partial \mu_{\mathrm{int}}}
                           + \omega \, (\mu_{\mathrm{cog}} - \mu_{\mathrm{int}}).
\end{align}

The coupling term $\omega$ controls the degree of emotional embodiment:
high $\omega$ produces strong bodily resonance (affective reactivity),
whereas low $\omega$ yields dissociation or purely cognitive inference.
This establishes a formal link between interoception, emotion,
and metacognition.

\subsection{Implementation Summary}

\begin{table}[h]
\centering
\caption{Physiological Extensions of Active Inference Agent}
\begin{tabular}{l l l}
\hline
Version & Key Innovation & Functional Role \\
\hline
v37.0 & Interoceptive generative model & Autonomic homeostasis \\
v38.0 & Energy metabolism field & Fatigue and motivational control \\
v39.0 & Heart–brain synchronization & Precision rhythm modulation \\
v40.0 & Interoceptive–cognitive coupling & Emotion and embodiment \\
\hline
\end{tabular}
\end{table}

\subsection{Summary}

The physiological active inference framework unifies metabolic regulation,
affective response, and cognitive control as aspects of a single
variational process.
Energy, heart rate, and internal states become precision variables
that dynamically shape confidence, motivation, and awareness.
This extension provides a principled route for modeling
emotion, stress, and bodily awareness within predictive coding networks.

% ==============================================================
\section{Social and Emotional Active Inference (v41–v45)}
% ==============================================================

\subsection{Rationale}

Versions v41–v45 expand the unified active inference agent into
the social–emotional domain.
Whereas v37–v40 linked cognition to internal physiology,
the new models describe how agents infer and regulate
\textit{other minds}—their intentions, beliefs, and emotions—
through shared predictive structures.
Social emotions such as empathy, trust, guilt, or envy
emerge as the natural consequence of minimizing
\textit{interpersonal free energy}.

\subsection{Interpersonal Generative Model}

Each agent $A_i$ maintains beliefs not only about its own latent states
$s^{(i)}$, but also about another agent’s inferred states $\hat{s}^{(j)}$.
The joint generative model factorizes as:

\[
p(o^{(i)}, s^{(i)}, \hat{s}^{(j)}) =
p(o^{(i)}|s^{(i)}) \,
p(s^{(i)}|\hat{s}^{(j)}) \,
p(\hat{s}^{(j)}),
\]
where $p(s^{(i)}|\hat{s}^{(j)})$ encodes
the agent’s \emph{social prior}—its beliefs about
how the other’s intentions affect its own state.
This coupling term allows agents to infer hidden causes behind
another’s actions, forming a minimal theory of mind.

\subsection{Interpersonal Free Energy}

The total free energy of agent $A_i$ is given by:

\[
F^{(i)}_{\mathrm{social}} =
F^{(i)}_{\mathrm{self}}
+ \xi \, \mathbb{E}_{q(\hat{s}^{(j)})}
\big[\| \hat{s}^{(j)} - s^{(i)} \|^2\big],
\]
where $\xi$ represents the precision of social coupling
(empathy gain).
Minimizing $F^{(i)}_{\mathrm{social}}$
aligns internal states across agents,
producing emotional resonance or empathy.

When $\xi$ is high, agents synchronize their affective trajectories;
when $\xi$ is low or asymmetric, misalignment occurs,
producing emotional isolation or misunderstanding.

\subsection{Emotional Valence and Precision}

Version v42 formalizes emotional valence as the sign of precision-weighted
prediction error in the interoceptive–social field:

\[
v^{(i)}_t = - \text{sgn}\!\big( \frac{\partial F^{(i)}_{\mathrm{social}}}{\partial \xi} \big),
\]
such that positive valence corresponds to successful alignment,
and negative valence reflects unexpected social feedback.
This unifies emotional appraisal with the same precision-based logic
that governs dopaminergic and interoceptive regulation.

\subsection{Empathy and Mirror Inference}

Version v43 introduces a mirror-inference field that reconstructs
another agent’s internal states from observed actions:

\[
\hat{s}^{(j)}_t = f_{\mathrm{mirror}}(o^{(j)}_t, a^{(j)}_t)
= W_{\mathrm{mirror}} \, [o^{(j)}_t, a^{(j)}_t],
\]
analogous to mirror-neuron processing.
Agents thus maintain internal simulations of others’ affective dynamics,
allowing prediction of cooperative or antagonistic intentions.

\subsection{Trust and Deception}

In v44, agents infer the reliability of another’s communication channel
via a \textit{trust precision} variable $\Pi_{\mathrm{trust}}$:

\[
\Pi_{\mathrm{trust}}^{(i,j)} =
\sigma \!\left(- D_{\mathrm{KL}}(q_i(s)|q_j(s)) \right),
\]
where $\sigma(\cdot)$ is the logistic function.
High divergence between beliefs reduces trust precision,
increasing defensive or competitive policy selection.
This mechanism captures both social learning and deception detection
as precision-weighting processes.

\subsection{Collective Emotion and Group Synchrony (v45)}

Multiple socially coupled agents form a group-level
affective field defined by the mean emotional energy:

\[
E_{\mathrm{group}} = \frac{1}{N}
\sum_{i=1}^{N} v^{(i)}_t \, \Pi_{\mathrm{trust}}^{(i)}.
\]

The temporal derivative $\dot{E}_{\mathrm{group}}$
reflects the collective affective momentum of the network.
When group precision is high and interoceptive priors are aligned,
collective synchrony (e.g., shared excitement or distress) emerges.
This provides a formal account of social contagion,
herd behavior, and group-level emotional coherence.

\subsection{Hierarchical Summary}

\begin{table}[h]
\centering
\caption{Social–Emotional Extensions of Active Inference Agent}
\begin{tabular}{l l l}
\hline
Version & Key Innovation & Functional Role \\
\hline
v41.0 & Interpersonal generative model & Mutual belief inference \\
v42.0 & Precision-based valence & Affective alignment \\
v43.0 & Mirror inference field & Empathy and action understanding \\
v44.0 & Trust precision variable & Reliability and deception inference \\
v45.0 & Group emotional field & Collective synchrony \\
\hline
\end{tabular}
\end{table}

\subsection{Summary}

The social–emotional expansion integrates interpersonal inference,
empathy, and collective affect into the unified variational framework.
Every interaction—alignment, competition, or misunderstanding—
emerges from reciprocal precision modulation across agents.
Thus, emotions are not added variables but intrinsic signals of
belief synchronization within a shared generative model.


% ==============================================================
\section{Cultural and Semantic Active Inference (v46–v50)}
% ==============================================================

\subsection{Rationale}

Versions v46–v50 elevate the active inference architecture
from social coordination (v41–v45) to shared symbolic
and cultural inference.
Here, groups of agents communicate and align through
language-like codes, shared narratives, and cultural priors.
Culture is treated as a distributed generative model that
minimizes collective uncertainty over meaning, norms, and goals.
This framework formalizes how semantic alignment, convention,
and ideology emerge as stable attractors in a network
of predictive agents.

\subsection{Semantic Generative Model}

Each agent $A_i$ encodes not only sensorimotor and emotional states
but also \emph{semantic variables} $s^{(\mathrm{sem})}$ that link
concepts, symbols, and meanings:

\[
p(o^{(i)}, s^{(i)}, s^{(\mathrm{sem})}) =
p(o^{(i)}|s^{(i)}) \,
p(s^{(i)}|s^{(\mathrm{sem})}) \,
p(s^{(\mathrm{sem})}).
\]

The semantic layer provides abstract priors $p(s^{(\mathrm{sem})})$
representing linguistic or cultural knowledge.
Language thus acts as a high-level generative model
that constrains perception, emotion, and social action.

\subsection{Cultural Priors and Normative Attractors}

In v47, each community of agents $\mathcal{C} = \{A_1, A_2, \dots\}$
shares a set of cultural priors $C_{\mathrm{cult}}$
encoding collective expectations about behavior and value.
These priors act as attractors in the belief space:

\[
p(s^{(\mathrm{sem})}|C_{\mathrm{cult}}) \propto
\exp\!\big(-\beta_{\mathrm{cult}} \|s^{(\mathrm{sem})} - \bar{s}_{\mathrm{group}}\|^2\big),
\]
where $\beta_{\mathrm{cult}}$ denotes cultural precision,
and $\bar{s}_{\mathrm{group}}$ is the group-level semantic centroid.
Cultural conformity and deviation (innovation) are then
two ends of a precision-controlled trade-off between stability and novelty.

\subsection{Language as Precision Transmission}

Version v48 introduces linguistic communication as precision exchange.
A message $m_t$ conveys the sender’s posterior mean
$\mathbb{E}[s^{(\mathrm{sem})}]$ weighted by confidence $\Pi_{\mathrm{msg}}$:

\[
m_t = \Pi_{\mathrm{msg}} \, \mathbb{E}_{q(s^{(\mathrm{sem})})}[s^{(\mathrm{sem})}],
\quad
\Pi_{\mathrm{msg}} = \frac{1}{1 + H[q(s^{(\mathrm{sem})})]}.
\]

Upon receipt, the listener updates its posterior according to
precision-weighted belief fusion:

\[
q_{\mathrm{new}}(s) \propto
q_{\mathrm{old}}(s)^{1-\Pi_{\mathrm{msg}}}
p(m_t|s)^{\Pi_{\mathrm{msg}}}.
\]

This formulation explains why clarity and confidence in communication
depend on both epistemic precision and shared priors.

\subsection{Norm Inference and Cultural Prediction Error (v49)}

Each agent tracks cultural prediction error as the discrepancy between
expected and observed social meaning:

\[
\varepsilon_{\mathrm{cult}} =
\big\| o^{(\mathrm{soc})}_t - g_{\mathrm{cult}}(s^{(\mathrm{sem})}_t) \big\|^2.
\]

A high $\varepsilon_{\mathrm{cult}}$ indicates norm violation or innovation,
which propagates as increased epistemic drive
and social uncertainty.
Agents with higher cognitive flexibility reduce $\varepsilon_{\mathrm{cult}}$
faster, modeling cultural adaptation and learning.

\subsection{Semantic Synchrony and Collective Meaning (v50)}

At the population level, shared meaning emerges as
synchronization in the semantic manifold:

\[
S_{\mathrm{group}} = \frac{1}{N}
\sum_i \mathbb{E}_{q_i(s^{(\mathrm{sem})})}[s^{(\mathrm{sem})}],
\quad
H_{\mathrm{group}} = -\sum_i H[q_i(s^{(\mathrm{sem})})].
\]

A high $H_{\mathrm{group}}$ corresponds to cultural diversity,
whereas low entropy (strong alignment) reflects consolidation
of a collective worldview.
Temporal evolution of $S_{\mathrm{group}}$
models cultural drift, convergence, or polarization
as emergent phenomena under variational optimization.

\subsection{Hierarchical Summary}

\begin{table}[h]
\centering
\caption{Cultural–Semantic Extensions of Active Inference Agent}
\begin{tabular}{l l l}
\hline
Version & Key Innovation & Functional Role \\
\hline
v46.0 & Semantic generative model & Conceptual meaning inference \\
v47.0 & Cultural priors & Normative stabilization \\
v48.0 & Language as precision transmission & Communication and belief fusion \\
v49.0 & Cultural prediction error & Norm learning and innovation \\
v50.0 & Semantic synchrony & Collective meaning formation \\
\hline
\end{tabular}
\end{table}

\subsection{Summary}

The cultural–semantic expansion generalizes active inference
from interpersonal interaction to symbolic and collective cognition.
Language, norms, and shared worldviews emerge as
high-level priors that constrain perception, emotion, and behavior.
Under this framework, cultural evolution itself
is an optimization of free energy across minds:
minimizing uncertainty about what is true, valued, and shared.


% ==============================================================
\section{Metacognitive and Reflective Active Inference (v51–v55)}
% ==============================================================

\subsection{Rationale}

Versions v51–v55 introduce a meta-level of inference:
agents not only infer external and internal states,
but also maintain beliefs about the \emph{reliability} of their own beliefs.
This “inference about inference” provides a formal mechanism for
self-monitoring, confidence evaluation, and reflective awareness.
Metacognitive processes emerge when an agent models the
uncertainty and consequences of its own epistemic actions.

\subsection{Metacognitive Generative Model}

The agent now includes an additional latent variable
$s^{(\mathrm{meta})}$ representing beliefs about belief quality.
The extended hierarchical model is:

\[
p(o, s^{(1)}, s^{(2)}, s^{(\mathrm{meta})})
= p(o|s^{(1)}) \, p(s^{(1)}|s^{(2)}) \, p(s^{(2)}|s^{(\mathrm{meta})}) \, p(s^{(\mathrm{meta})}).
\]

Here $s^{(1)}$ and $s^{(2)}$ correspond to sensory and cognitive states,
while $s^{(\mathrm{meta})}$ generates predictions about
the reliability, precision, and confidence of lower levels.
This model explicitly encodes epistemic uncertainty and
allows the system to adapt its precision hierarchy.

\subsection{Confidence as Meta-Precision}

Version v51 defines confidence as a meta-precision parameter
$\Pi_{\mathrm{meta}}$, representing the expected reliability of
the agent’s own posterior beliefs:

\[
\Pi_{\mathrm{meta}} = \sigma\!\big( -H[q(s^{(2)})] \big),
\]
where $H[q]$ is the entropy of the second-level posterior.
High entropy (uncertainty) lowers meta-precision,
signaling doubt or hesitation,
while low entropy increases confidence and decisiveness.
This mechanism unifies metacognitive confidence with
the precision control already used at perceptual and social levels.

\subsection{Reflective Update and Meta-Prediction Error (v52)}

Agents now compute a meta-prediction error,
defined as the discrepancy between predicted and realized inference quality:

\[
\varepsilon_{\mathrm{meta}} =
|\widehat{\Pi}_{\mathrm{meta}} - \Pi_{\mathrm{meta}}|.
\]
Large meta-errors trigger introspective updates,
modifying the parameters of the internal model
(e.g., learning rate, attention gain, or memory decay).
This provides a formal definition of reflection:
a second-order correction to the inference process itself.

\subsection{Counterfactual and Alternative Model Simulation (v53)}

Version v53 introduces counterfactual simulation:
agents can generate alternative generative models
$\tilde{M}_k$ and evaluate their free energy
without physically acting:

\[
\tilde{F}_k = \mathbb{E}_{q_k}[\ln q_k - \ln p_k(o, s)].
\]
The policy posterior becomes a mixture over real and imagined models,
enabling mental simulation, planning, and creative reasoning.
This parallels human reflection and internal deliberation:
testing beliefs against simulated possibilities.

\subsection{Meta-Learning and Precision Adaptation (v54)}

The agent’s learning rates $\alpha_t$ and precision gains $\Pi_t$
are now subject to meta-level adaptation:

\[
\alpha_{t+1} = \alpha_t + \eta_\alpha (\varepsilon_{\mathrm{meta}} - \bar{\varepsilon}),
\qquad
\Pi_{t+1} = \Pi_t + \eta_\Pi (\Pi_{\mathrm{meta}} - \bar{\Pi}),
\]
where $\eta_\alpha$ and $\eta_\Pi$ are meta-learning rates.
This allows the system to tune its own inference parameters
based on historical performance—an intrinsic form of self-optimization.

\subsection{Self-Model and Introspective Loop (v55)}

In v55, the agent explicitly models itself as part of its generative world.
The self-model $M_{\mathrm{self}}$ predicts the consequences
of the agent’s own actions and thoughts:

\[
p(o, a | M_{\mathrm{self}}) = p(o|a, s) \, p(a|M_{\mathrm{self}}).
\]
Discrepancies between predicted and experienced internal states
produce introspective free energy,
which drives reconfiguration of the self-model.

This introduces recursive self-awareness:
the agent not only acts but also represents itself as acting—
a necessary precondition for metacognition, agency, and identity.

\subsection{Hierarchical Summary}

\begin{table}[h]
\centering
\caption{Metacognitive Extensions of Active Inference Agent}
\begin{tabular}{l l l}
\hline
Version & Key Innovation & Functional Role \\
\hline
v51.0 & Meta-precision variable & Confidence monitoring \\
v52.0 & Meta-prediction error & Reflective correction \\
v53.0 & Counterfactual simulation & Planning and creativity \\
v54.0 & Meta-learning of parameters & Self-optimization \\
v55.0 & Self-model and introspection & Identity and agency \\
\hline
\end{tabular}
\end{table}

\subsection{Summary}

The metacognitive layer transforms the agent
from an adaptive inference system into a self-reflective one.
It monitors uncertainty, simulates alternatives,
and adapts its own inference parameters.
This extension provides a principled account of introspection,
deliberation, and self-awareness as hierarchical free-energy minimization
applied recursively to the process of inference itself.


% ============================================================
\section{Implementation: Connectome-Constrained Spiking Zebrafish Brain (v2)}
% ============================================================

\subsection{Overview}

We implemented the theoretical framework as a large-scale spiking neural network
embedded in a 2D prey--predator virtual environment.
The model comprises \textbf{7,369 Izhikevich spiking neurons} across 22 brain modules,
constrained by the MPIN-Atlas (MapZebrain) zebrafish connectome.
All inter-region projections follow anatomically validated pathways
(Retina $\to$ Tectum $\to$ Thalamus $\to$ Pallium $\to$ Subpallium),
with distance-dependent connection probabilities derived from atlas centroids.

\subsection{Neuron and Synapse Models}

Each neuron follows the Izhikevich two-variable model with midpoint integration:
\begin{align}
v_{t+1} &= v_t + \frac{\Delta t}{2}\left(0.04v_t^2 + 5v_t + 140 - u_t + I_t\right), \\
u_{t+1} &= u_t + \frac{\Delta t}{2}\,a\,(b\,v_t - u_t),
\end{align}
with spike reset $v \to c$, $u \to u + d$ when $v \geq 30$\,mV.
Seven cell types are used: regular spiking (RS), intrinsic bursting (IB),
fast spiking (FS), chattering (CH), low-threshold spiking (LTS),
thalamocortical (TC), and medium spiny neuron (MSN).

Synaptic transmission uses conductance-based dynamics with three receptor types:
\begin{equation}
I_{\mathrm{syn}} = g \cdot (E_{\mathrm{rev}} - V_{\mathrm{post}}),
\end{equation}
where $g = \bar{g} \cdot W \cdot s$ and the gating variable $s$ decays exponentially
with time constants $\tau_{\mathrm{AMPA}} = 5$\,ms, $\tau_{\mathrm{NMDA}} = 80$\,ms,
$\tau_{\mathrm{GABA_A}} = 10$\,ms.
NMDA currents include voltage-dependent Mg$^{2+}$ block.

\subsection{Neural Architecture}

All layers use 75\% excitatory / 25\% inhibitory (E/I) balance with four
intra-layer synapse populations ($E \to E$, $E \to I$, $I \to E$, $I \to I$).
Fifty Izhikevich substeps per behavioral step provide 1\,ms temporal resolution.

\begin{table*}[h!]
\centering
\caption{Neural architecture: layer sizes and cell types}
\renewcommand{\arraystretch}{1.15}
\begin{tabular}{l r r l l}
\hline
\textbf{Region} & \textbf{E neurons} & \textbf{I neurons} & \textbf{Cell types} & \textbf{Atlas region} \\
\hline
Retina (4 RGC types)     & \multicolumn{2}{c}{rate-coded}   & ON/OFF/Loom/DS    & R \\
Tectum SFGS-b/d           & 1800 & 600  & CH, FS   & TeO \\
Tectum SGC/SO             & 600  & 200  & IB/RS, FS & TeO \\
Thalamus (TC + TRN)       & 300  & 80   & LTS, FS  & Th \\
Pallium-S (superficial)   & 1200 & 400  & RS, FS   & P (Dl) \\
Pallium-D (deep)          & 600  & 200  & IB, FS   & P (Dm) \\
Basal ganglia (D1/D2/GPi) & 760  & ---  & MSN      & SP \\
Amygdala (LA/CeA/ITC)     & 20   & 10   & RS/IB/FS & Dm \\
Reticulospinal            & 42   & ---  & Named    & aRF/imRF \\
Place cells               & 400  & ---  & $\theta$-gated & Dl \\
Goal selector (WTA)       & 4    & ---  & RS       & --- \\
\hline
\textbf{Total}            & \multicolumn{2}{c}{\textbf{7,369}} & 7 types & 8 atlas regions \\
\hline
\end{tabular}
\end{table*}

\subsection{Connectome Constraints}

Inter-region connectivity follows the MPIN-Atlas (MapZebrain) with 36 brain regions
and 17 validated projection pathways.
Connection probability between regions $i$ and $j$ is:
\begin{equation}
P(i \to j) = s_{ij} \cdot p_{\mathrm{base}},
\end{equation}
where $s_{ij} \in [0, 1]$ is the anatomical projection strength from published
tract-tracing data, and $p_{\mathrm{base}} = 0.3$ sets maximum connectivity density.

Key projections include:
Retina $\to$ Tectum ($s = 1.0$, contralateral),
Tectum $\to$ Thalamus ($s = 0.6$),
Thalamus $\to$ Pallium ($s = 0.7$, reciprocal),
Pallium $\to$ Subpallium ($s = 0.6$),
Tectum $\to$ Reticular formation ($s = 0.7$, escape pathway).

\subsection{Predictive Coding in the Pallium}

The pallium implements a two-layer predictive coding hierarchy:
\begin{itemize}
\item \textbf{Pallium-S} (superficial): receives thalamic relay, computes somatic rate $r_S$.
\item \textbf{Pallium-D} (deep): goal/intent representation, projects to basal ganglia.
\item \textbf{Feedback}: $W_{\mathrm{FB}}$ projects Pallium-D $\to$ Pallium-S (apical prediction).
\end{itemize}

Prediction error is computed as the apical--somatic mismatch:
\begin{equation}
\varepsilon_t = W_{\mathrm{FB}} \cdot r_D - r_S,
\end{equation}
and drives anti-Hebbian feedback learning:
\begin{equation}
\Delta W_{\mathrm{FB}} = -\eta \cdot \pi \cdot \varepsilon_t \cdot r_D^\top.
\end{equation}

Free energy is:
\begin{equation}
F_t = \|\varepsilon_t\|^2.
\end{equation}

\subsection{Active Inference and Goal Selection}

Goal selection uses a hybrid architecture:
\begin{enumerate}
\item \textbf{Analytic EFE}: Expected free energy per goal policy computed from
    food evidence, threat evidence, energy state, allostatic errors,
    place cell memory, and world model predictions.
\item \textbf{Spiking WTA}: Four-neuron winner-take-all attractor network receiving
    Pallium-D population code and EFE bias. Self-excitatory ($w = 4.0$) and
    mutually inhibitory ($w = -2.0$) connections create bistable attractor dynamics.
\item \textbf{Fallback}: When WTA confidence $< 0.4$, analytic EFE determines goal.
\end{enumerate}

The WTA network implements policy selection as an attractor settling process,
providing a neural mechanism for active inference goal commitment
with hysteresis (persistence) and sharp transitions.

\subsection{Neuromodulation}

Four neuromodulatory axes operate as scalar signals updated per behavioral step:
\begin{itemize}
\item \textbf{Dopamine (DA)}: Reward prediction error via TD(0),
    $\delta_t = r_t - V_t$, modulates D1/D2 balance in basal ganglia.
\item \textbf{Noradrenaline (NA)}: Arousal signal driven by amygdala threat.
\item \textbf{Serotonin (5-HT)}: Patience/habituation, suppresses flee during PATROL states.
\item \textbf{Acetylcholine (ACh)}: Attention/plasticity gate, modulates pallial precision.
\end{itemize}

\subsection{Plasticity}

Three-factor spike-timing-dependent plasticity (STDP) with eligibility traces:
\begin{align}
e_{ij}(t) &= \lambda \, e_{ij}(t-1) + A^+ s_j^{\mathrm{post}} \bar{s}_i^{\mathrm{pre}} - A^- \bar{s}_j^{\mathrm{post}} s_i^{\mathrm{pre}}, \\
\Delta W_{ij} &= \eta \cdot \mathrm{DA} \cdot \mathrm{ACh} \cdot e_{ij},
\end{align}
where $\bar{s}$ denotes the eligibility trace filtered spike train,
and consolidation is gated by both dopamine (reward relevance) and
acetylcholine (attentional state).

Episodic fear conditioning in the amygdala uses Hebbian LTP:
$\Delta W_{\mathrm{LA} \to \mathrm{CeA}} = 0.05 \cdot s^{\mathrm{CeA}} \cdot (s^{\mathrm{LA}})^\top$
triggered by near-death events (proximity $> 0.9$ and predator facing fish).

\subsection{Autonomous Perception}

The model operates without access to ground-truth environment state,
relying on three sensory modalities:

\paragraph{Retinal input.}
A geometric sensory bridge projects environmental objects onto 800-pixel
bilateral retinal arrays (400 intensity + 400 type channel per eye)
with a 200$^\circ$ binocular field of view (rear blind spot reduced
from 220$^\circ$ to 160$^\circ$ vs.\ prior 140$^\circ$ FoV).
A central binocular overlap zone (40$^\circ$) is encoded by both eyes
at reduced intensity.
Four retinal ganglion cell types (ON-sustained, OFF-transient, looming,
direction-selective) encode the visual scene.
Food effective radius is 12~px (up from 8~px), modeling motion shimmer
detection by ON-sustained RGCs.

\paragraph{Predator tracking.}
Predator position is estimated via a Kalman filter driven by retinal observations:
\begin{align}
\hat{x}_{t|t-1} &= \hat{x}_{t-1} + \hat{v}_{t-1}, \\
K_t &= P_t / (P_t + R), \\
\hat{x}_t &= \hat{x}_{t|t-1} + K_t (z_t - \hat{x}_{t|t-1}),
\end{align}
where $z_t$ is estimated from retinal enemy pixel count and lateral bias.
Object permanence is maintained when the predator leaves the visual field.
Flee direction uses \textbf{direct retinal escape with symmetry-breaking}:
the fish turns away from the eye with more enemy pixels; when enemy pixels
are equal in both eyes (predator centered), a random direction is chosen
(Mauthner cell lateralization), preventing decision paralysis.
A temporal EMA filter ($\alpha=0.3$) smooths the enemy lateral bias
to reduce directional jitter.
This bypasses the Kalman position estimate which proved unreliable
for directional control.

\paragraph{Lateral line mechanoreception.}
A simulated lateral line provides non-visual proximity detection within
$\sim$150\,px (body-length range) with 15\,px Gaussian noise.
This supplements retinal tracking for predators approaching from
the rear hemisphere (outside 200$^\circ$ FoV).

\paragraph{Rear-threat inference.}
When amygdala arousal is elevated ($\alpha > 0.15$) but enemy pixels
drop to zero, the system infers the predator is behind the fish and
maintains the last known flee direction with exponential decay
($\tau = 0.9$), implementing a form of working memory for threat location.

\paragraph{Self-localization.}
Fish position is estimated from place cell population vector:
$\hat{p} = \sum_i w_i c_i / \sum_i w_i$,
where $w_i$ is the firing rate and $c_i$ the centroid of place cell $i$.
Heading is proprioceptive (vestibular).

\subsection{Allostatic Regulation}

Three interoceptive variables drive homeostatic goal biases:
\begin{itemize}
\item \textbf{Hunger}: $h_t = 1 - E_t / 100$, biases toward FORAGE.
\item \textbf{Fatigue}: accumulates with speed, biases toward EXPLORE (rest).
\item \textbf{Stress}: driven by predator proximity, biases toward FLEE.
\end{itemize}

Each variable computes an allostatic error (predicted $-$ setpoint)
that additively modulates the expected free energy of each goal policy,
implementing interoceptive inference within the active inference framework.

\subsection{Results}

\subsubsection{Decision Quality}

The model was evaluated on five structured decision scenarios
(Table~\ref{tab:scenarios}) comparing v2 (spiking, connectome-constrained)
against v1 (rate-coded, unconstrained):

\begin{table}[h!]
\centering
\caption{Decision scenario scores (max 100). Three model configurations:
v1 (rate-coded, ground-truth access), v2-GT (spiking, 2.3K neurons, ground-truth),
v2-auto (spiking, 7.4K neurons, retinal + lateral line only).}
\label{tab:scenarios}
\begin{tabular}{l c c c}
\hline
\textbf{Scenario} & \textbf{v1 (rate)} & \textbf{v2-GT (2.3K)} & \textbf{v2-auto (7.4K)} \\
\hline
A: Safe vs Risky food   & 100 & 100 & 80 \\
B: Occluded food        & 85  & 76  & 66 \\
C: Predator charge      & 80  & 100 & 88 \\
D: Starvation dilemma   & 100 & 100 & 93 \\
E: Detour               & 79  & 95  & 55 \\
\hline
\textbf{Overall}        & 89  & \textbf{94} & \textbf{77} \\
\hline
\end{tabular}
\end{table}

Three key results emerge.
First, with ground-truth sensory access, the spiking v2 model \textbf{exceeded}
the rate-coded v1 baseline (94 vs.\ 89), demonstrating that Izhikevich
dynamics, conductance-based synapses, and E/I balance do not degrade
decision quality and may improve it through richer temporal dynamics.

Second, the fully autonomous v2-auto model achieves 77/100 using only
retinal vision (200$^\circ$ FoV), lateral line mechanoreception, and
place cell self-localization---with no ground-truth environment reads.
The 17-point gap from v2-GT was reduced from 22 points through five
sensory improvements: FoV widened from 140$^\circ$ to 200$^\circ$,
food detection radius increased from 8~px to 12~px (motion shimmer),
symmetry-breaking flee implementing Mauthner cell lateralization,
binocular overlap zone in the central 40$^\circ$, and temporal EMA
bearing filter ($\alpha=0.3$).
All five scenarios now achieve survival; the remaining gap is attributable
to retinal distance estimation being less precise than ground-truth,
reducing food interception efficiency (Scenarios A and E).

Third, the predator charge scenario (C) shows a notable \textbf{v2 advantage}:
the spiking model's flee response is faster (first flee at step 0) than v1
(which required several steps to build up threat evidence), likely because
the Izhikevich IB neurons in the amygdala CeA produce bursting responses
to sudden threat onset.

\subsubsection{Ablation Study}

Systematic removal of individual modules reveals their contributions:

\begin{table}[h!]
\centering
\caption{Ablation study: module contributions}
\begin{tabular}{l c c}
\hline
\textbf{Condition} & \textbf{Score} & \textbf{$\Delta$} \\
\hline
Full v2 (GT mode)  & 94 & --- \\
No dopamine        & 92 & $-2$ \\
No amygdala        & 92 & $-2$ \\
No allostasis      & 85 & $-9$ \\
No world model     & 92 & $-2$ \\
\hline
\end{tabular}
\end{table}

Allostatic regulation is the most critical module ($-9$ points),
confirming that interoceptive inference drives the starvation--safety trade-off
central to adaptive behavior.

\subsubsection{STDP Training}

Reward-modulated STDP training over 200 episodes (4-stage curriculum)
was performed on the full 7,369-neuron model using an NVIDIA RTX~A6000 GPU
(14.2 hours wall-clock time). The curriculum progressively introduces
environmental complexity:

\begin{table}[h!]
\centering
\caption{STDP curriculum training results (200 episodes, 7.4K neurons)}
\begin{tabular}{l c c c c}
\hline
\textbf{Stage} & \textbf{Steps} & \textbf{Food/ep} & \textbf{Survival} & \textbf{Flee \%} \\
\hline
1: Forage easy  & 300 & 1.70 & 100\% & 31\% \\
2: Forage hard  & 400 & 0.98 & 100\% & 42\% \\
3: Survive easy & 500 & 1.68 & 86\%  & 52\% \\
4: Survive hard & 600 & 1.00 & 72\%  & 48\% \\
\hline
\textbf{Overall} & --- & 1.34 & \textbf{90\%} & 45\% \\
\hline
\end{tabular}
\end{table}

The survival rate decreasing from Stage~3 (86\%) to Stage~4 (72\%)
reflects the difficulty of sparse food combined with an active predator.
The flee fraction increases monotonically across stages, indicating
appropriate threat-response calibration.

\subsubsection{Layer Activity}

After tuning subthreshold tonic drive ($I_{\mathrm{tonic}}^E = -2.0$\,pA,
$I_{\mathrm{tonic}}^I = 0.5$--$1.0$\,pA) and normalized input currents
(target 3.0\,pA mean, clamped at 6.0\,pA max), all spiking layers show
stimulus-dependent, sparse activity consistent with biological recordings:

\begin{table}[h!]
\centering
\caption{Mean firing rates per layer (spikes/ms) at 7.4K scale.
Biological target: tectum 0.001--0.05, pallium 0.001--0.02.}
\begin{tabular}{l r c l}
\hline
\textbf{Layer} & \textbf{Neurons} & \textbf{Mean rate} & \textbf{Regime} \\
\hline
Tectum SFGS-b E & 900  & 0.003 & Sparse, input-driven \\
Tectum SFGS-d E & 900  & 0.007 & Sparse, OFF-transient \\
Tectum SGC E    & 300  & 0.002 & Looming-selective \\
Thalamus TC     & 300  & 0.033 & Tonic relay \\
Thalamus TRN    & 80   & 0.009 & Gating \\
Pallium-S E     & 1200 & 0.005 & Moderate \\
Pallium-S I     & 400  & active & Tonic ($I_\mathrm{tonic}=2.0$\,pA) \\
Pallium-D E     & 600  & 0.002 & Goal-dependent \\
BG D1           & 400  & 0.040 & DA-modulated \\
BG gate         & ---  & 0.43  & Action gating \\
Amygdala CeA    & 20   & 0.007 & Threat-proportional \\
Place cells     & 400  & 0.006 & Theta-gated sparse \\
\hline
\end{tabular}
\end{table}

\subsection{Discussion}

The v2 implementation demonstrates that a connectome-constrained spiking network
with biologically realistic neuron models, conductance-based synapses, and
local plasticity rules can reproduce the core behavioral repertoire of larval zebrafish:
prey capture, predator avoidance, energy management, and goal-directed exploration.

\paragraph{Cost of biological realism (reduced gap).}
The gap between v2-GT (94/100) and v2-auto was reduced from 22 to 17 points
(72/100 $\to$ 77/100) through five sensory improvements.
The primary bottleneck was the 140$^\circ$ FoV leaving a 220$^\circ$ rear
blind spot; widening to 200$^\circ$ (160$^\circ$ blind spot) combined with
symmetry-breaking Mauthner flee and lateral line proximity detection
eliminated predator-capture death in Scenarios C and D.
The residual 17-point gap is attributable to retinal pixel-count-to-distance
mapping being less precise than ground-truth coordinates for food approach
(Scenarios A and E) and obstacle-occluded food detection (Scenarios B and E).
Both limitations are biologically realistic---real zebrafish compensate through
scanning saccades and schooling behavior.

\paragraph{Advantages of spiking dynamics.}
Despite the perceptual cost, the spiking model shows specific advantages.
In the predator charge scenario (C), v2 achieves faster flee initiation
(step~0 vs.\ step~3 for v1) because IB-type amygdala neurons produce
burst responses to sudden threat onset, providing a biologically faithful
implementation of the Mauthner-mediated fast escape reflex.

\paragraph{Computational cost.}
The 7,369-neuron model with 50 Izhikevich substeps per behavioral step
and conductance-based synapses requires $\sim$1.4 million operations per step.
STDP training (200 episodes) requires $\sim$14 hours on an RTX~A6000.
Real-time simulation at 30\,fps is not achievable at this scale;
future optimization through vectorized substep batching and GPU kernel fusion
could reduce this by an estimated 5--10$\times$.

\paragraph{Future directions.}
Three extensions would address the remaining 17-point performance gap:
(i) active scanning saccades to further reduce the 160$^\circ$ rear blind spot
(FoV already widened to 200$^\circ$),
(ii) integration of the full MapZebrain fiber tract data for
region-to-region connectivity,
and (iii) multi-agent schooling behavior, where conspecific alarm
signals provide indirect predator detection from angles not
directly visible to the focal fish.

The framework provides a platform for testing active inference theories
in a biologically grounded whole-brain model with experimentally measurable
predictions about neural dynamics, neuromodulatory effects, and behavioral
phenotypes.


% ============================================================
\section{Acknowledgments}

Additional information can be given in the template, such as to not include funder information in the acknowledgments section.

\nocite{*} % This command displays all refs in the bib file. PLEASE DELETE IT BEFORE YOU SUBMIT YOUR MANUSCRIPT!
\bibliography{references}

%%%%%%%%%%%%%%%%%%%%%%%%%%%%%%%%%%%%%%%%%%%%%%%%%%%%%%%%%%%%
%%% APPENDICES
%%%%%%%%%%%%%%%%%%%%%%%%%%%%%%%%%%%%%%%%%%%%%%%%%%%%%%%%%%%%

\appendix

\end{document}



